% Back matter plan: statements required by Nature Portfolio journals. Updated 17 Nov 2026.

\section*{Declarations}

\role{Required statements; data and code availability also carry the release argument of the paper.}

\structure
\noindent{\small
\begin{tabularx}{\linewidth}{@{}>{\raggedright\arraybackslash}p{0.22\linewidth}>{\raggedright\arraybackslash}X>{\raggedright\arraybackslash}p{0.26\linewidth}@{}}
\toprule
Statement & What goes in & Open decision \\ \midrule
Data availability & Benchmark (cases, policies, annotations) with a DOI; MIMIC-IV via PhysioNet DUA
  & Test answers with a canary string, or behind an evaluation server \\
Code availability & Harness, reward, training, evaluation; Zenodo DOI at the paper commit, available
  to reviewers at submission & --- \\
Acknowledgements & Reviewers not authors; CERN compute & --- \\
Funding & Nature Portfolio format & Grant numbers \\
Author contributions & CRediT roles & Author list \\
Competing interests & Statement & --- \\
Ethics & Synthetic cases, no patient data; consent or exemption for review and rating study; MIMIC-IV
  DUA & Check with supervisors \\
Reporting & Nature Portfolio Reporting Summary (exact $P$, $n$, CIs); MI-CLAIM-GEN and TRIPOD-LLM
  checklists & --- \\
\bottomrule
\end{tabularx}}

\status{Decide the test-split release before submission; it also affects the public repository.}
